%% problem_id: S047-ASNSP_2013_5_12_4_1001_0-prop-3.4
%% source_id: S047 (NUMDAM; Annali della Scuola Normale Superiore di Pisa, Cl. Sci. (5))
%% source_article: Marina Ghisi, Massimo Gobbino, "Higher order Glaeser inequalities and optimal
%%   regularity of roots of real functions", Ann. Sc. Norm. Super. Pisa Cl. Sci. (5) XII (2013),
%%   1001--1021.
%% source_url: https://www.numdam.org/item/ASNSP_2013_5_12_4_1001_0.pdf
%% representation: PDF.  The text layer is UNRELIABLE, so the statement and the proof below were
%%   transcribed by RENDERING the pages (pdftoppm -r 150 -png) and READING THE IMAGES, then checking
%%   every display against the page.
%% statement_locator: printed page 1011 (= PDF page 11), Proposition 3.4, display (3.5)
%% proof_locator: printed pages 1011--1012 (= PDF pages 11--12), from "Proof." on p. 1011 to the
%%   closing square on p. 1012
%% difficulty_level: H3
%% why THIS item and not Theorem 2.1: Theorem 2.1's entire proof is one paragraph reducing to this
%%   proposition ("the assumptions of Proposition 3.4 are satisfied for every d > 0", plus
%%   Hoeld_a(v^(k),(x0-d,x0+d)) <= Hoeld_a(v^(k),R), then d -> +infinity).  So the substance lives
%%   HERE and, per plan 2.2 second criterion, this is the item worth collecting.
%% dedup_note: 同篇的 Theorem 2.1 是它的一段化归；按 §2.2 第二判据收本条而非那条。两条不是同一命题。
%% status: complete (statement + author proof verbatim + reason + locators)
%%
%% The statement and proof are the AUTHORS' text; nothing is paraphrased, supplemented or guessed.
%% ---------------------------------------------------------------------------
\documentclass[11pt]{article}
\usepackage{amsmath,amssymb,amsthm}
\newtheorem{proposition}{Proposition}
\newcommand{\Hold}{\operatorname{H\"old}}
\begin{document}

\section*{Proposition 3.4 (Higher order Glaeser inequality in an interval)}

\textbf{Source.} M.~Ghisi, M.~Gobbino, \emph{Higher order Glaeser inequalities and optimal
regularity of roots of real functions}, Ann.\ Sc.\ Norm.\ Super.\ Pisa Cl.\ Sci.\ (5)
\textbf{XII} (2013), 1001--1021; printed page 1011.

\begin{proposition}
Let $k$ be a positive integer, and let $\alpha \in (0,1]$, $x_0 \in \mathbb{R}$, and $d>0$ be
real numbers.

Let $v \in C^{k}((x_0-d,\,x_0+d))$ be a function such that $v(x)$ and $v'(x)$ do not change their
sign, namely either $v(x)\ge 0$ or $v(x)\le 0$ for every $x \in (x_0-d,\,x_0+d)$, and the same for
$v'(x)$. Let us assume that the $k$-th derivative $v^{(k)}(x)$ is $\alpha$-H\"older continuous in
$(x_0-d,\,x_0+d)$.

Then there exists a constant $C(k)$, depending only upon $k$, such that
\begin{equation}
|v'(x_0)|^{k+\alpha}
\;\le\; C(k)\cdot |v(x_0)|^{k+\alpha-1}\cdot
\max\left\{ \Hold_\alpha\!\left(v^{(k)},\,(x_0-d,\,x_0+d)\right),\;
\frac{|v'(x_0)|}{d^{\,k+\alpha-1}} \right\}. \tag{3.5}
\end{equation}
\end{proposition}

\subsection*{Proof (authors' text)}

Up to symmetries we can assume, without loss of generality, that $v(x)\ge 0$ and $v'(x)\le 0$ for
every $x \in (x_0-d,\,x_0+d)$. If $v'(x_0)=0$, then there is nothing to prove. If $v(x_0)=0$, then
it is easy to see that also $v'(x_0)=0$, and once again there is nothing to prove. So from now on
we assume that $v(x_0)>0$ and $v'(x_0)<0$.

Let $\delta_k$ be the constant of Lemma 3.3. We distinguish two cases.

\medskip\noindent\textbf{Case 1.} Let us assume that
\[
\frac{|v'(x_0)|}{v(x_0)}\, d \;\le\; \delta_k .
\]
This is equivalent to saying that $|v'(x_0)| \le \delta_k\,d^{-1}v(x_0)$, and hence
\[
|v'(x_0)|^{k+\alpha} \;=\; |v'(x_0)|\cdot|v'(x_0)|^{k+\alpha-1}
\;\le\; |v'(x_0)|\cdot\left(\frac{\delta_k}{d}\right)^{k+\alpha-1}[v(x_0)]^{k+\alpha-1},
\]
which implies (3.5) in this first case, provided that $C(k)\ge \delta_k^{\,k}$.

\medskip\noindent\textbf{Case 2.} Let us assume now that
\begin{equation}
\frac{|v'(x_0)|}{v(x_0)}\, d \;>\; \delta_k . \tag{3.6}
\end{equation}
Let us set for simplicity $H := \Hold_\alpha\!\left(v^{(k)},\,(x_0-d,\,x_0+d)\right)$, and let us
write Taylor's expansion of $v(x)$ of order $k$. For every $\delta \in [0,d]$ we obtain that
\[
v(x_0+\delta) \;=\; \sum_{j=0}^{k-1}\frac{v^{(j)}(x_0)}{j!}\,\delta^{j}
\;+\; \frac{v^{(k)}(x_0+\xi)}{k!}\,\delta^{k}
\]
for some $\xi \in (0,\delta)$. Since $v^{(k)}$ is $\alpha$-H\"older continuous we have that
\[
\left|v^{(k)}(x_0+\xi)-v^{(k)}(x_0)\right| \;\le\; H\,\delta^{\alpha}.
\]
Therefore our assumption that $v(x)\ge 0$ for every $x \in (x_0-d,\,x_0+d)$ implies that
\[
0 \;\le\; v(x_0+\delta) \;\le\; \sum_{j=0}^{k}\frac{v^{(j)}(x_0)}{j!}\,\delta^{j}
\;+\; \frac{H}{k!}\,\delta^{k+\alpha} \qquad \forall \delta \in [0,d].
\]
In an analogous way we obtain that
\[
0 \;\ge\; v'(x_0+\delta) \;\ge\; \sum_{j=0}^{k-1}\frac{v^{(j+1)}(x_0)}{j!}\,\delta^{j}
\;-\; \frac{H}{(k-1)!}\,\delta^{k+\alpha-1} \qquad \forall \delta \in [0,d].
\]
Let us write both inequalities with $\delta := x\cdot v(x_0)\cdot|v'(x_0)|^{-1}$. We obtain that
the two inequalities
\[
0 \le |v(x_0)|\left(\sum_{j=0}^{k}\frac{v^{(j)}(x_0)}{j!}\cdot\frac{[v(x_0)]^{j-1}}{|v'(x_0)|^{j}}
\cdot x^{j} + \frac{H}{k!}\cdot\frac{[v(x_0)]^{k+\alpha-1}}{|v'(x_0)|^{k+\alpha}}\cdot x^{k+\alpha}\right),
\]
\[
0 \ge |v'(x_0)|\left(\sum_{j=0}^{k-1}\frac{v^{(j+1)}(x_0)}{j!}\cdot\frac{[v(x_0)]^{j}}{|v'(x_0)|^{j+1}}
\cdot x^{j} - \frac{H}{(k-1)!}\cdot\frac{[v(x_0)]^{k+\alpha-1}}{|v'(x_0)|^{k+\alpha}}
\cdot x^{k+\alpha-1}\right)
\]
hold true for every
\[
0 \le x < \frac{|v'(x_0)|}{v(x_0)}\, d .
\]
Due to (3.6), the upper bound is larger than $\delta_k$, which proves that the two inequalities hold
true at least for every $x \in [0,\delta_k]$. Thus if we set
\begin{equation}
P(x) := \sum_{j=0}^{k}\frac{v^{(j)}(x_0)}{j!}\cdot\frac{[v(x_0)]^{j-1}}{|v'(x_0)|^{j}}\cdot x^{j},
\qquad
A := \frac{H}{k!}\cdot\frac{[v(x_0)]^{k+\alpha-1}}{|v'(x_0)|^{k+\alpha}}, \tag{3.7}
\end{equation}
we have that $P(0)=1$, $P'(0)=-1$, and the assumptions in the left-hand side of (3.2) are satisfied.

Therefore from (3.2) it follows that $A \ge \varepsilon_k$, which implies (3.5) also in this second
case, provided that $C(k) \ge (\varepsilon_k k!)^{-1}$. \hfill$\square$

\end{document}
